\section{Praktik: Membangun Aplikasi Web Sederhana (CRUD Dasar)}

Section ini menggabungkan seluruh materi Bab 11 ke dalam satu studi kasus utuh: \textbf{Aplikasi Daftar Tugas} (\textit{to-do list}) dengan fitur CRUD, session login, dan validasi input. Aplikasi terdiri dari beberapa file terpisah yang mengikuti prinsip \textit{separation of concerns}: file konfigurasi database, file koneksi, halaman login, dashboard CRUD, dan logout. Arsitektur ini mencerminkan pola umum pengembangan aplikasi web PHP di dunia nyata \cite{phpManual}.

\begin{figure}[ht]
\centering
\begin{tikzpicture}[
  box/.style={draw, rounded corners, minimum width=2.5cm, minimum height=1cm, align=center, font=\footnotesize, fill=#1},
  arrow/.style={-Stealth, thick},
  node distance=0.8cm
]
  \node[box=blue!15] (login) {login.php\\{\tiny Form login}};
  \node[box=green!15, right=2cm of login] (dash) {dashboard.php\\{\tiny Daftar \& form tugas}};
  \node[box=orange!15, below=1.5cm of dash] (proses) {proses.php\\{\tiny INSERT/UPDATE/DELETE}};
  \node[box=red!15, right=2cm of dash] (logout) {logout.php\\{\tiny Hapus session}};
  \node[box=yellow!20, below=1.5cm of login] (koneksi) {koneksi.php\\{\tiny mysqli connect}};
  \node[box=purple!15, below=of koneksi] (db) {MySQL\\{\tiny tabel: users, tugas}};

  \draw[arrow] (login) -- node[above, font=\scriptsize]{session} (dash);
  \draw[arrow] (dash) -- (proses);
  \draw[arrow] (proses) -- node[right, font=\scriptsize]{redirect} (dash);
  \draw[arrow] (dash) -- node[above, font=\scriptsize]{logout} (logout);
  \draw[arrow] (logout) -- node[above, font=\scriptsize]{redirect} (login);
  \draw[arrow] (koneksi) -- (db);
  \draw[arrow, dashed] (login) -- (koneksi);
  \draw[arrow, dashed] (dash) -- (koneksi);
  \draw[arrow, dashed] (proses) -- (koneksi);
\end{tikzpicture}
\caption{Arsitektur aplikasi daftar tugas (CRUD + session)}
\end{figure}

\begin{contoh}
\textbf{Studi Kasus: Aplikasi Daftar Tugas dengan Login dan CRUD}

Aplikasi ini memiliki alur sebagai berikut: (1) pengguna login melalui form --- data diverifikasi terhadap database menggunakan \code{password\_verify()}; (2) jika valid, session dimulai dan pengguna diarahkan ke dashboard; (3) di dashboard, pengguna dapat melihat daftar tugas (SELECT), menambah tugas baru (INSERT), mengedit tugas (UPDATE), dan menghapus tugas (DELETE); (4) semua operasi database menggunakan prepared statement; (5) logout menghancurkan session.
\end{contoh}

Langkah pertama adalah menyiapkan file koneksi database yang akan di-include oleh semua halaman yang membutuhkan akses database. File ini memisahkan konfigurasi koneksi agar tidak perlu diulang di setiap halaman dan memudahkan perubahan konfigurasi di satu tempat.

\begin{webcode}[language=PHP, caption={koneksi.php --- File koneksi database}]
<?php
// koneksi.php
$host     = 'localhost';
$username = 'root';
$password = '';
$database = 'crud_tugas';

$conn = new mysqli($host, $username, $password, $database);
if ($conn->connect_error) {
  die("Koneksi gagal: " . $conn->connect_error);
}
$conn->set_charset("utf8mb4");
?>
\end{webcode}

Berikut adalah skrip SQL untuk membuat tabel \code{users} dan \code{tugas} yang digunakan oleh aplikasi. Tabel \code{users} menyimpan data login, sedangkan tabel \code{tugas} menyimpan item to-do milik pengguna.

\begin{webcode}[language=SQL, caption={Skrip SQL untuk membuat tabel}]
CREATE DATABASE IF NOT EXISTS crud_tugas;
USE crud_tugas;

CREATE TABLE users (
  id INT AUTO_INCREMENT PRIMARY KEY,
  username VARCHAR(50) UNIQUE NOT NULL,
  password VARCHAR(255) NOT NULL
);

CREATE TABLE tugas (
  id INT AUTO_INCREMENT PRIMARY KEY,
  user_id INT NOT NULL,
  judul VARCHAR(200) NOT NULL,
  deskripsi TEXT,
  selesai TINYINT(1) DEFAULT 0,
  created_at TIMESTAMP DEFAULT CURRENT_TIMESTAMP,
  FOREIGN KEY (user_id) REFERENCES users(id)
);

-- Insert user demo (password: rahasia123)
INSERT INTO users (username, password) VALUES
  ('admin', '$2y$10$abcdefghijklmnopqrstuuABCDEFGHIJKLMNOPQRSTUVWXYZ12');
-- Gunakan password_hash('rahasia123', PASSWORD_DEFAULT) untuk hash asli
\end{webcode}

Halaman login memproses form POST, memverifikasi kredensial terhadap database, dan memulai session jika berhasil. Perhatikan penggunaan \code{password\_verify()} dan \code{session\_regenerate\_id()}.

\begin{webcode}[language=PHP, caption={login.php --- Halaman login dengan session}]
<?php
// login.php
session_start();
require_once 'koneksi.php';

$error = '';
if ($_SERVER['REQUEST_METHOD'] === 'POST') {
  $user = trim($_POST['username'] ?? '');
  $pass = $_POST['password'] ?? '';

  $stmt = $conn->prepare("SELECT id, username, password FROM users WHERE username = ?");
  $stmt->bind_param("s", $user);
  $stmt->execute();
  $result = $stmt->get_result();

  if ($row = $result->fetch_assoc()) {
    if (password_verify($pass, $row['password'])) {
      session_regenerate_id(true);
      $_SESSION['user_id']  = $row['id'];
      $_SESSION['username'] = $row['username'];
      header("Location: dashboard.php");
      exit;
    }
  }
  $error = "Username atau password salah.";
  $stmt->close();
}
?>
<!DOCTYPE html>
<html lang="id">
<head><title>Login</title></head>
<body>
  <h2>Login Daftar Tugas</h2>
  <?php if ($error): ?>
    <p style="color:red;"><?= htmlspecialchars($error) ?></p>
  <?php endif; ?>
  <form method="POST">
    <label>Username: <input type="text" name="username" required></label><br>
    <label>Password: <input type="password" name="password" required></label><br>
    <button type="submit">Login</button>
  </form>
</body>
</html>
\end{webcode}

Dashboard menampilkan daftar tugas pengguna dan menyediakan form untuk menambah atau mengedit tugas. Halaman ini juga mengirimkan request DELETE dan UPDATE melalui form POST ke \code{proses.php}.

\begin{webcode}[language=PHP, caption={dashboard.php --- CRUD daftar tugas (Read + form)}]
<?php
// dashboard.php
session_start();
if (!isset($_SESSION['user_id'])) {
  header("Location: login.php");
  exit;
}
require_once 'koneksi.php';

$userId = $_SESSION['user_id'];

// READ - ambil semua tugas milik user
$stmt = $conn->prepare("SELECT * FROM tugas WHERE user_id = ? ORDER BY created_at DESC");
$stmt->bind_param("i", $userId);
$stmt->execute();
$tugas = $stmt->get_result()->fetch_all(MYSQLI_ASSOC);
$stmt->close();

// Cek apakah dalam mode edit
$editData = null;
if (isset($_GET['edit'])) {
  $stmtEdit = $conn->prepare("SELECT * FROM tugas WHERE id = ? AND user_id = ?");
  $stmtEdit->bind_param("ii", $_GET['edit'], $userId);
  $stmtEdit->execute();
  $editData = $stmtEdit->get_result()->fetch_assoc();
  $stmtEdit->close();
}
?>
<!DOCTYPE html>
<html lang="id">
<head><title>Dashboard</title></head>
<body>
  <h2>Halo, <?= htmlspecialchars($_SESSION['username']) ?>!
      <a href="logout.php">Logout</a></h2>

  <h3><?= $editData ? 'Edit Tugas' : 'Tambah Tugas' ?></h3>
  <form method="POST" action="proses.php">
    <input type="hidden" name="aksi" value="<?= $editData ? 'update' : 'create' ?>">
    <?php if ($editData): ?>
      <input type="hidden" name="id" value="<?= $editData['id'] ?>">
    <?php endif; ?>
    <label>Judul: <input type="text" name="judul" required
      value="<?= htmlspecialchars($editData['judul'] ?? '') ?>"></label><br>
    <label>Deskripsi: <textarea name="deskripsi"><?=
      htmlspecialchars($editData['deskripsi'] ?? '') ?></textarea></label><br>
    <button type="submit"><?= $editData ? 'Simpan' : 'Tambah' ?></button>
  </form>

  <h3>Daftar Tugas</h3>
  <table border="1" cellpadding="5">
    <tr><th>Judul</th><th>Deskripsi</th><th>Tanggal</th><th>Aksi</th></tr>
    <?php foreach ($tugas as $t): ?>
    <tr>
      <td><?= htmlspecialchars($t['judul']) ?></td>
      <td><?= htmlspecialchars($t['deskripsi']) ?></td>
      <td><?= $t['created_at'] ?></td>
      <td>
        <a href="dashboard.php?edit=<?= $t['id'] ?>">Edit</a>
        <form method="POST" action="proses.php" style="display:inline;">
          <input type="hidden" name="aksi" value="delete">
          <input type="hidden" name="id" value="<?= $t['id'] ?>">
          <button type="submit" onclick="return confirm('Hapus?')">Hapus</button>
        </form>
      </td>
    </tr>
    <?php endforeach; ?>
  </table>
</body>
</html>
\end{webcode}

File \code{proses.php} menangani operasi CREATE, UPDATE, dan DELETE berdasarkan parameter \code{aksi} dari form. Semua operasi menggunakan prepared statement untuk keamanan.

\begin{webcode}[language=PHP, caption={proses.php --- Handler CREATE, UPDATE, DELETE}]
<?php
// proses.php
session_start();
if (!isset($_SESSION['user_id'])) {
  header("Location: login.php");
  exit;
}
require_once 'koneksi.php';

$userId = $_SESSION['user_id'];
$aksi   = $_POST['aksi'] ?? '';

switch ($aksi) {
  case 'create':
    $judul     = trim($_POST['judul'] ?? '');
    $deskripsi = trim($_POST['deskripsi'] ?? '');
    if (!empty($judul)) {
      $stmt = $conn->prepare(
        "INSERT INTO tugas (user_id, judul, deskripsi) VALUES (?, ?, ?)"
      );
      $stmt->bind_param("iss", $userId, $judul, $deskripsi);
      $stmt->execute();
      $stmt->close();
    }
    break;

  case 'update':
    $id        = intval($_POST['id'] ?? 0);
    $judul     = trim($_POST['judul'] ?? '');
    $deskripsi = trim($_POST['deskripsi'] ?? '');
    if ($id > 0 && !empty($judul)) {
      $stmt = $conn->prepare(
        "UPDATE tugas SET judul = ?, deskripsi = ? WHERE id = ? AND user_id = ?"
      );
      $stmt->bind_param("ssii", $judul, $deskripsi, $id, $userId);
      $stmt->execute();
      $stmt->close();
    }
    break;

  case 'delete':
    $id = intval($_POST['id'] ?? 0);
    if ($id > 0) {
      $stmt = $conn->prepare(
        "DELETE FROM tugas WHERE id = ? AND user_id = ?"
      );
      $stmt->bind_param("ii", $id, $userId);
      $stmt->execute();
      $stmt->close();
    }
    break;
}

$conn->close();
header("Location: dashboard.php");
exit;
?>
\end{webcode}

Terakhir, file \code{logout.php} menghancurkan session dan mengarahkan kembali ke halaman login.

\begin{webcode}[language=PHP, caption={logout.php --- Menghancurkan session}]
<?php
// logout.php
session_start();
$_SESSION = [];
session_destroy();
setcookie(session_name(), '', time() - 3600, '/');
header("Location: login.php");
exit;
?>
\end{webcode}

Studi kasus ini mengintegrasikan konsep-konsep utama Bab 11: OOP (class MySQLi), session (autentikasi login), cookie (session ID), validasi input (\code{trim()}, cek kosong), sanitasi output (\code{htmlspecialchars()}), prepared statement (mencegah SQL injection), dan arsitektur file yang terpisah. Sebagai latihan lanjutan, mahasiswa dapat menambahkan fitur: (1) registrasi pengguna baru dengan \code{password\_hash()}; (2) CSRF token pada setiap form; (3) pagination untuk daftar tugas; (4) fitur pencarian tugas; (5) upload lampiran file pada tugas \cite{phpRightWay}.

